\section{En Python}
\label{sec:python}

Chaque outil de ce chapitre se vérifie en quelques lignes de NumPy. Les sorties indiquées en commentaire sont celles obtenues ; ces programmes sont regroupés dans le notebook \texttt{TP0.ipynb}.

\subsection{Matrice de covariance et valeurs propres}
\begin{python}
import numpy as np

# Matrice de covariance et valeurs propres (sections 1)
X = np.array([[1, 2], [2, 2], [3, 5]])        # 3 individus, 2 variables
Xc = X - X.mean(axis=0)                        # tableau centré
S = Xc.T @ Xc / len(X)                         # covariance empirique
print(S.round(3))                              # [[0.667 1.] [1. 2.]]
lam, U = np.linalg.eigh(S)                     # valeurs propres croissantes
print(lam.round(4))                            # [0.1315 2.5352]
print(U[:, -1].round(3))                       # premier axe de l'ACP
\end{python}
On retrouve la matrice $S$ de l'exemple de la section~\ref{sec:algebre}. Le vecteur propre de la plus grande valeur propre donne la direction de plus grande dispersion des trois points : c'est le premier axe de l'ACP.

\subsection{Droite des moindres carrés et descente de gradient}
\begin{python}
import numpy as np

# Droite des moindres carrés (section 2)
x = np.array([1, 2, 3]); y = np.array([2, 2, 5])
b = np.mean((x - x.mean()) * (y - y.mean())) / np.var(x)   # s_xy / s_x²
a = y.mean() - b * x.mean()
print(a, b)                                    # 0.0 1.5
print(np.polyfit(x, y, deg=1).round(4))        # [1.5 0.] : même droite

# Descente de gradient sur L(θ) = (θ - 3)²
theta, eta = 0.0, 0.25
for k in range(6):
    theta = theta - eta * 2 * (theta - 3)
    print(round(theta, 4), end=" ")            # 1.5 2.25 2.625 2.8125 ...
\end{python}
Les formules de la proposition~\ref{prop:droite} donnent la même droite que \texttt{np.polyfit}, et la descente de gradient s'approche de $3$ comme prévu par la proposition~\ref{prop:gd}.

\subsection{Loi des grands nombres et théorème central limite}
\begin{python}
import numpy as np
rng = np.random.default_rng(2)

# Loi des grands nombres et TCL (section 3)
p, n, R = 0.3, 400, 20_000                     # R sondages de n personnes
props = rng.binomial(n, p, size=R) / n         # R proportions observées
print(props.mean().round(4), props.var().round(6), p * (1 - p) / n)
ecart = 1.96 * np.sqrt(p * (1 - p) / n)
print(np.mean(np.abs(props - p) <= ecart).round(3))   # proche de 0.95
\end{python}
La variance des proportions observées vaut bien $p(1-p)/n$, et l'écart à $p$ reste sous $1{,}96\sqrt{p(1-p)/n}$ dans environ $95\,\%$ des sondages simulés ($94{,}3\,\%$ ici).

\subsection{Moyennes conditionnelles}
\begin{python}
import numpy as np, pandas as pd
rng = np.random.default_rng(3)

# Moyennes conditionnelles par groupe (section 4)
n = 10_000
milieu = rng.choice(["urbain", "rural"], size=n, p=[0.4, 0.6])
enfants = np.where(milieu == "urbain", rng.poisson(2.5, n), rng.poisson(4.5, n))
df = pd.DataFrame({"milieu": milieu, "enfants": enfants})

print(df.groupby("milieu")["enfants"].mean().round(2))   # ≈ 4.5 et 2.5
part = df["milieu"].value_counts(normalize=True)
moy = df.groupby("milieu")["enfants"].mean()
print(round((part * moy).sum(), 4), round(df["enfants"].mean(), 4))  # espérance totale
\end{python}
Les moyennes par groupe estiment $\E[Y\mid X=x]$, et leur moyenne pondérée redonne la moyenne générale (proposition~\ref{prop:tour}).
